% !TeX program = xelatex
\documentclass[12pt,a4paper]{article}
\usepackage{fontspec}
\setmainfont{Times New Roman}
\usepackage[vietnamese]{babel}
\usepackage{amsmath,amssymb}
\usepackage[margin=2.3cm]{geometry}
\usepackage{enumitem}
\setlist[enumerate]{leftmargin=*,itemsep=0.4em}
\setlength{\parindent}{0pt}
\begin{document}
\begin{center}
{\Large\bfseries ĐỀ TOÁN TUYỂN SINH VÀO LỚP 10 THPT HÀ NỘI -- NĂM HỌC 2024--2025}\\[2pt]
{\large\bfseries PHẦN ĐẠI SỐ VÀ TOÁN ỨNG DỤNG}\\
Môn Toán -- Ngày thi: 09/06/2024
\end{center}

\textit{Các câu hình học chứng minh được lược bỏ.}

\section*{Câu I. Biểu thức đại số}
Cho hai biểu thức
\[
A=\frac{x}{\sqrt{x}-3},\qquad
B=\frac{2x-3}{x-3\sqrt{x}}-\frac1{\sqrt{x}},
\quad x>0,\ x\ne9.
\]
\begin{enumerate}
\item Tính giá trị của $A$ khi $x=16$.
\item Chứng minh rằng $B=\dfrac{2\sqrt{x}-1}{\sqrt{x}-3}$.
\item Tìm tất cả giá trị của $x$ để $A-B<0$.
\end{enumerate}

\section*{Câu II. Bài toán thực tế}
\begin{enumerate}
\item Để chở 15 tấn thiết bị, một đội vận chuyển dự định dùng xe tải loại nhỏ. Đội chuyển sang chỉ dùng xe tải loại lớn, vì vậy số xe giảm 2 chiếc so với dự định và mỗi xe tải loại lớn chở nhiều hơn mỗi xe loại nhỏ 2 tấn. Biết mỗi xe cùng loại chở số tấn hàng bằng nhau. Hỏi đội sử dụng bao nhiêu xe tải loại lớn?
\item Một bình đựng nước hình trụ có bán kính đáy 4 cm và chiều cao 25 cm. Tính diện tích xung quanh của bình (lấy $\pi\approx3{,}14$).
\end{enumerate}

\section*{Câu III. Hệ phương trình và hàm số}
\begin{enumerate}
\item Giải hệ phương trình
\[
\begin{cases}
\sqrt{3x+1}+2y=4,\\
3\sqrt{3x+1}-y=5.
\end{cases}
\]
\item Trong mặt phẳng tọa độ $Oxy$, cho parabol $(P):y=x^2$ và đường thẳng $(d):y=(m-2)x+5$.
\begin{enumerate}
\item Chứng minh $(d)$ luôn cắt $(P)$ tại hai điểm phân biệt.
\item Gọi $x_1,x_2$ là hoành độ các giao điểm. Tìm tất cả giá trị của $m$ để $x_1+5x_2=0$.
\end{enumerate}
\end{enumerate}

\section*{Câu V. Bài toán tối ưu}
Với các số thực dương $x,y$ thỏa mãn $x+y+xy=3$, tìm giá trị nhỏ nhất của
\[
P=\frac{3}{x+y}-xy.
\]
\end{document}
